% Generated draft: requires mathematical and source review.
\documentclass[11pt]{article}
\usepackage[margin=25mm]{geometry}
\usepackage{fontspec}
\setmainfont{Latin Modern Roman}
\usepackage{amsmath,amssymb,mathtools,unicode-math}
\setmathfont{Latin Modern Math}
\usepackage{enumitem,needspace,xurl,hyperref}
\setlength{\parindent}{0pt}
\setlength{\parskip}{5pt}
\setlength{\emergencystretch}{3em}
\newcommand{\N}{\mathbb{N}}
\newcommand{\Z}{\mathbb{Z}}
\newcommand{\Q}{\mathbb{Q}}
\newcommand{\R}{\mathbb{R}}
\newcommand{\C}{\mathbb{C}}
\newcommand{\F}{\mathbb{F}}
\newcommand{\A}{\mathbb{A}}
\newcommand{\PP}{\mathbb{P}}
\newcommand{\E}{\mathbb{E}}
\newcommand{\eps}{\varepsilon}
\newcommand{\dd}{\,\mathrm{d}}
\begin{document}
\section*{UnsolvedMath ID: 30004120}
\textbf{Source problem code:} catalogue entry
\subsection{Definitions and mathematical statement}
Let \(n\geq 3\) be an odd integer, let \([n]=\{1,\ldots,n\}\) represent the people, and let \(K_n\) be the complete simple graph on \([n]\). Under the usual graph interpretation, a round table seating \(\ell\geq 3\) people determines a cycle \(C_\ell\), whose edges are the unordered pairs of adjacent people. A meal seating all \(n\) people therefore determines a spanning \(2\)-regular simple graph, called a \(2\)-factor. Its components describe the tables. Requiring every unordered pair of distinct people to be adjacent at exactly one meal means that the edge sets of the meal graphs partition \(E(K_n)\). Each meal contributes \(n\) edges, so the prescribed number of meals is \((n-1)/2\); in particular, this formulation requires odd \(n\).

\textbf{Conditional precise formulation of the intended Oberwolfach question.} Interpret ``variable table sizes'' as allowing an arbitrary prescribed list of table sizes, with that same list used at every meal. Is there an integer \(N\geq 3\) such that, for every odd integer \(n\geq N\), every integer \(m\geq 1\), and every list of integers \(\ell_1,\ldots,\ell_m\geq 3\) satisfying \(\sum_{i=1}^{m}\ell_i=n\), there exist spanning subgraphs \(H_1,\ldots,H_{(n-1)/2}\) of \(K_n\) such that
\[
H_j\cong F:=\bigsqcup_{i=1}^{m}C_{\ell_i}
\quad\text{for every }1\leq j\leq (n-1)/2,
\qquad
E(K_n)=\bigsqcup_{j=1}^{(n-1)/2}E(H_j)?
\]
Here \(\cong\) means graph isomorphism, and the edge union is disjoint. The threshold \(N\) is uniform over all prescribed lists of table sizes.

\textbf{Unresolved source qualifications.} The supplied seating sentence does not explicitly prescribe a table-size list, require its repetition at every meal, exclude tables of size one or two, or specify whether the sufficiently-large threshold is uniform over table configurations. Thus the displayed statement is a conditional reconstruction using the supplied identification with the Oberwolfach problem. If table sizes may change between meals, the condition \(H_j\cong F\) would not follow from the sentence; if the table configuration is merely existentially chosen, its quantifiers would also differ. No even-order variant is specified.

The supplied registry labels the record \texttt{partially\_solved}, while its statement-verification note says that the official report answers the sufficiently-large-order question affirmatively and that this extraction is a proved result rather than an open problem. Its statement-recovery status is \texttt{unrecoverable}, and no clean statement is supplied. These are supplied curation claims, not independently verified mathematical conclusions. The separately mentioned full classification of all two-factor instances is not substituted for the stated sufficiently-large-order question.
\subsection{Short English statement}
Under the intended interpretation, prescribe any collection of round tables, each seating at least three people, with an odd total of \(n\) seats. The question asks whether, for all sufficiently large such \(n\), the same table sizes can be used for \((n-1)/2\) meals so that every pair of people sits next to each other exactly once. The source leaves these table-size quantifiers implicit. Its supplied verification note reports that the sufficiently-large-order question has already been answered affirmatively; this draft does not independently verify that claim.
\subsection{Sources}
\begin{itemize}
\item Supplied local UnsolvedMath registry record, permanent ID \(30004120\), title \emph{Oberwolfach Seating Problem}, exact source code \texttt{catalogue entry}, in the \emph{Oberwolfach Reports Open Problems} collection. The proposed year is recorded as 2019; no individual proposer is supplied.
\item Supplied source citation: \emph{Combinatorics, Probability and Computing (2020)}, Oberwolfach Reports, DOI 10.4171/owr/2019/19, pp.~1113--1189. Source URL and supplied literature-source link labelled ``Oberwolfach Report 19/2019'': \url{https://doi.org/10.4171/owr/2019/19}. The supplied background describes these reports as records of mathematical workshops and their open-problem discussions.
\item Supplied literature source: Glock--Joos--Kim--Kuhn--Osthus, \emph{Resolution of the Oberwolfach problem}, \url{https://doi.org/10.4171/JEMS/1060}.
\item Literature assessment dated exactly \texttt{2026-08-22}: the supplied curation attributes the sufficiently-large admissible-order result to Glock, Joos, Kim, Kuhn, and Osthus, and describes the full classification of all two-factor instances as remaining open apart from known exceptions and finite computations. This is supplied dated curation, not independently verified. The supplied background expressly requests independent verification of time-sensitive status and cited claims; no such verification is asserted here.
\item The supplied statement-verification note reports that the official report explicitly answers the extracted sufficiently-large-order question affirmatively, retains the original record for provenance, and promotes no separate clean statement. This source-specific qualification accompanies the registry's supplied \texttt{partially\_solved} label.
\end{itemize}
\end{document}
